\chapter{Discrete Joint Distributions}
\addtolength{\topmargin}{-0.9pt}

\section{Basic Concepts}
\index{joint distribution}
\index{joint PMF}

\begin{definition}[Joint Probability Mass Function]\index{joint PMF|see{joint probability mass function}}
Let $X$ and $Y$ be discrete random variables. The \emph{joint probability mass function (joint PMF)} of $X$ and $Y$ is defined as:
\[
p_{X,Y}(x, y) = P(X = x \text{ and } Y = y), \quad \text{for all } x, y
\]

The joint PMF satisfies:
\begin{enumerate}
    \item $p_{X,Y}(x, y) \geq 0$ for all $x, y$
    \item $\sum_{x} \sum_{y} p_{X,Y}(x, y) = 1$
\end{enumerate}
\end{definition}

\begin{definition}[Marginal PMF]\index{marginal PMF}
The \emph{marginal PMF} of $X$ is:
\[
p_X(x) = \sum_{y} p_{X,Y}(x, y)
\]

Similarly, the marginal PMF of $Y$ is:
\[
p_Y(y) = \sum_{x} p_{X,Y}(x, y)
\]
\end{definition}

\begin{exercise}[DIY: Marginals Sum to 1]
Prove that if $p_{X,Y}(x, y)$ is a valid joint PMF, then:
\[
\sum_{x} p_X(x) = 1 \quad \text{and} \quad \sum_{y} p_Y(y) = 1
\]
\end{exercise}

\begin{proof}
\[
\sum_{x} p_X(x) = \sum_{x} \sum_{y} p_{X,Y}(x, y) = 1
\]
since the joint PMF sums to 1. Similarly:
\[
\sum_{y} p_Y(y) = \sum_{y} \sum_{x} p_{X,Y}(x, y) = 1
\]
\end{proof}

\begin{eg}[Simple Joint Distribution]
Consider two random variables $X$ and $Y$ with the following joint PMF:

\begin{center}
\renewcommand{\arraystretch}{1.5}
\begin{tabular}{c|ccc}
$p_{X,Y}(x,y)$ & $X=1$ & $X=2$ & $X=3$ \\ \hline
$Y=2$ & 0.1 & 0.2 & 0.1 \\
$Y=3$ & 0.15 & 0.15 & 0.1 \\
$Y=4$ & 0.05 & 0.1 & 0.05 \\
\end{tabular}
\end{center}

Find:
\begin{enumerate}[label=(\alph*)]
    \item The marginal PMFs of $X$ and $Y$
    \item $P(X = 2)$
    \item $P(Y \geq 3)$
    \item $P(X = 2, Y = 3)$
\end{enumerate}

\begin{solution}
\begin{enumerate}[label=(\alph*)]
    \item \textbf{Marginal PMF of $X$:}
    \[
    p_X(1) = 0.1 + 0.15 + 0.05 = 0.3, \quad p_X(2) = 0.2 + 0.15 + 0.1 = 0.45, \quad p_X(3) = 0.1 + 0.1 + 0.05 = 0.25
    \]
    
    \textbf{Marginal PMF of $Y$:}
    \[
    p_Y(2) = 0.1 + 0.2 + 0.1 = 0.4, \quad p_Y(3) = 0.15 + 0.15 + 0.1 = 0.4, \quad p_Y(4) = 0.05 + 0.1 + 0.05 = 0.2
    \]
    
    \item $P(X = 2) = p_X(2) = 0.45$
    
    \item $P(Y \geq 3) = p_Y(3) + p_Y(4) = 0.4 + 0.2 = 0.6$
    
    \item $P(X = 2, Y = 3) = p_{X,Y}(2, 3) = 0.15$
\end{enumerate}
\end{solution}
\end{eg}

\section{Conditional Distributions}
\index{conditional distribution}
\index{conditional PMF}

\begin{definition}[Conditional PMF]
The \emph{conditional PMF} of $Y$ given $X = x$ is:
\[
p_{Y|X}(y \mid x) = \frac{p_{X,Y}(x, y)}{p_X(x)}, \quad \text{provided } p_X(x) > 0
\]

Similarly, the conditional PMF of $X$ given $Y = y$ is:
\[
p_{X|Y}(x \mid y) = \frac{p_{X,Y}(x, y)}{p_Y(y)}, \quad \text{provided } p_Y(y) > 0
\]
\end{definition}

\begin{exercise}[DIY: Conditional PMF is Valid PMF]
Prove that for a fixed $x$ with $p_X(x) > 0$, the conditional PMF $p_{Y|X}(y \mid x)$ satisfies:
\begin{enumerate}[label=(\alph*)]
    \item $p_{Y|X}(y \mid x) \geq 0$ for all $y$
    \item $\sum_{y} p_{Y|X}(y \mid x) = 1$
\end{enumerate}
\end{exercise}

\begin{proof}
\begin{enumerate}[label=(\alph*)]
    \item Since $p_{X,Y}(x, y) \geq 0$ and $p_X(x) > 0$, we have $p_{Y|X}(y \mid x) = \frac{p_{X,Y}(x, y)}{p_X(x)} \geq 0$.
    
    \item 
    \[
    \sum_{y} p_{Y|X}(y \mid x) = \sum_{y} \frac{p_{X,Y}(x, y)}{p_X(x)} = \frac{1}{p_X(x)} \sum_{y} p_{X,Y}(x, y) = \frac{p_X(x)}{p_X(x)} = 1
    \]
\end{enumerate}
\end{proof}

\begin{eg}[Conditional Distributions]
Using the joint distribution from the previous example, find:
\begin{enumerate}[label=(\alph*)]
    \item $p_{Y|X}(y \mid 2)$ (the conditional PMF of $Y$ given $X = 2$)
    \item $p_{X|Y}(x \mid 3)$ (the conditional PMF of $X$ given $Y = 3$)
    \item $P(Y = 4 \mid X = 1)$
\end{enumerate}

\begin{solution}
\begin{enumerate}[label=(\alph*)]
    \item Given $X = 2$, we have $p_X(2) = 0.45$:
    \[
    p_{Y|X}(2 \mid 2) = \frac{0.2}{0.45} = \frac{4}{9}, \quad p_{Y|X}(3 \mid 2) = \frac{0.15}{0.45} = \frac{1}{3}, \quad p_{Y|X}(4 \mid 2) = \frac{0.1}{0.45} = \frac{2}{9}
    \]
    
    \item Given $Y = 3$, we have $p_Y(3) = 0.4$:
    \[
    p_{X|Y}(1 \mid 3) = \frac{0.15}{0.4} = \frac{3}{8}, \quad p_{X|Y}(2 \mid 3) = \frac{0.15}{0.4} = \frac{3}{8}, \quad p_{X|Y}(3 \mid 3) = \frac{0.1}{0.4} = \frac{1}{4}
    \]
    
    \item $P(Y = 4 \mid X = 1) = p_{Y|X}(4 \mid 1) = \frac{p_{X,Y}(1, 4)}{p_X(1)} = \frac{0.05}{0.3} = \frac{1}{6}$
\end{enumerate}
\end{solution}
\end{eg}

\section{Independence}
\index{independence!discrete random variables}

\begin{definition}[Independence]
Two discrete random variables $X$ and $Y$ are \emph{independent} if and only if:
\[
p_{X,Y}(x, y) = p_X(x) \cdot p_Y(y) \quad \text{for all } x, y
\]

Equivalently, $X$ and $Y$ are independent if:
\[
p_{Y|X}(y \mid x) = p_Y(y) \quad \text{for all } x, y \text{ with } p_X(x) > 0
\]
\end{definition}

\begin{exercise}[DIY: Equivalence of Independence Definitions]
Prove that the two definitions of independence are equivalent:
\begin{enumerate}[label=(\alph*)]
    \item If $p_{X,Y}(x, y) = p_X(x) \cdot p_Y(y)$ for all $x, y$, then $p_{Y|X}(y \mid x) = p_Y(y)$ for all $x, y$ with $p_X(x) > 0$.
    \item If $p_{Y|X}(y \mid x) = p_Y(y)$ for all $x, y$ with $p_X(x) > 0$, then $p_{X,Y}(x, y) = p_X(x) \cdot p_Y(y)$ for all $x, y$.
\end{enumerate}
\end{exercise}

\begin{proof}
\begin{enumerate}[label=(\alph*)]
    \item If $p_{X,Y}(x, y) = p_X(x) \cdot p_Y(y)$, then:
    \[
    p_{Y|X}(y \mid x) = \frac{p_{X,Y}(x, y)}{p_X(x)} = \frac{p_X(x) \cdot p_Y(y)}{p_X(x)} = p_Y(y)
    \]
    provided $p_X(x) > 0$.
    
    \item If $p_{Y|X}(y \mid x) = p_Y(y)$ for all $x, y$ with $p_X(x) > 0$, then:
    \[
    p_{X,Y}(x, y) = p_{Y|X}(y \mid x) \cdot p_X(x) = p_Y(y) \cdot p_X(x) = p_X(x) \cdot p_Y(y)
    \]
    For $x$ with $p_X(x) = 0$, we have $p_{X,Y}(x, y) = 0 = p_X(x) \cdot p_Y(y)$ as well.
\end{enumerate}
\end{proof}

\begin{eg}[Checking Independence]
Determine if $X$ and $Y$ are independent in the previous example.

\begin{solution}
We check if $p_{X,Y}(x, y) = p_X(x) \cdot p_Y(y)$ for all $x, y$.

For example, check $(x, y) = (1, 2)$:
\[
p_{X,Y}(1, 2) = 0.1
\]
\[
p_X(1) \cdot p_Y(2) = 0.3 \cdot 0.4 = 0.12
\]

Since $0.1 \neq 0.12$, $X$ and $Y$ are \textbf{not independent}.
\end{solution}
\end{eg}

\begin{eg}[Independent Random Variables]
Consider two independent random variables $X$ and $Y$ with:
\[
p_X(1) = 0.4, \quad p_X(2) = 0.6
\]
\[
p_Y(2) = 0.3, \quad p_Y(3) = 0.5, \quad p_Y(4) = 0.2
\]

Find the joint PMF $p_{X,Y}(x, y)$.

\begin{solution}
Since $X$ and $Y$ are independent:
\[
p_{X,Y}(x, y) = p_X(x) \cdot p_Y(y)
\]

\begin{center}
\renewcommand{\arraystretch}{1.5}
\begin{tabular}{c|cc}
$p_{X,Y}(x,y)$ & $X=1$ & $X=2$ \\ \hline
$Y=2$ & $0.4 \cdot 0.3 = 0.12$ & $0.6 \cdot 0.3 = 0.18$ \\
$Y=3$ & $0.4 \cdot 0.5 = 0.20$ & $0.6 \cdot 0.5 = 0.30$ \\
$Y=4$ & $0.4 \cdot 0.2 = 0.08$ & $0.6 \cdot 0.2 = 0.12$ \\
\end{tabular}
\end{center}
\end{solution}
\end{eg}

\section{Expectation and Functions of Two Random Variables}
\index{expectation!joint distributions}
\index{linearity of expectation}

\begin{exercise}[DIY: Linearity of Expectation]
Prove that for any constants $a, b, c$:
\[
E(aX + bY + c) = aE(X) + bE(Y) + c
\]

\textbf{Hint:} Use the definition $E[g(X, Y)] = \sum_{x} \sum_{y} g(x, y) \cdot p_{X,Y}(x, y)$.
\end{exercise}

\begin{proof}
\[
E(aX + bY + c) = \sum_{x} \sum_{y} (ax + by + c) \cdot p_{X,Y}(x, y)
\]
\[
= a \sum_{x} \sum_{y} x \cdot p_{X,Y}(x, y) + b \sum_{x} \sum_{y} y \cdot p_{X,Y}(x, y) + c \sum_{x} \sum_{y} p_{X,Y}(x, y)
\]
\[
= a \sum_{x} x \sum_{y} p_{X,Y}(x, y) + b \sum_{y} y \sum_{x} p_{X,Y}(x, y) + c \cdot 1
\]
\[
= a \sum_{x} x \cdot p_X(x) + b \sum_{y} y \cdot p_Y(y) + c = aE(X) + bE(Y) + c
\]

This holds \textbf{regardless of whether $X$ and $Y$ are independent}!
\end{proof}

\begin{eg}[Expectation Calculations]
Using the joint distribution from the first example, find:
\begin{enumerate}[label=(\alph*)]
    \item $E(X)$ and $E(Y)$
    \item $E(XY)$
    \item $E(2X + 3Y - 1)$
    \item $E(X^2)$
\end{enumerate}

\begin{solution}
\begin{enumerate}[label=(\alph*)]
    \item 
    \[
    E(X) = 1(0.3) + 2(0.45) + 3(0.25) = 0.3 + 0.9 + 0.75 = 1.95
    \]
    \[
    E(Y) = 2(0.4) + 3(0.4) + 4(0.2) = 0.8 + 1.2 + 0.8 = 2.8
    \]
    
    \item 
    \[
    E(XY) = \sum_{x} \sum_{y} xy \cdot p_{X,Y}(x, y)
    \]
    \[
    = 1 \cdot 2(0.1) + 1 \cdot 3(0.15) + 1 \cdot 4(0.05) + 2 \cdot 2(0.2) + 2 \cdot 3(0.15)
    \]
    \[
    + 2 \cdot 4(0.1) + 3 \cdot 2(0.1) + 3 \cdot 3(0.1) + 3 \cdot 4(0.05)
    \]
    \[
    = 0.2 + 0.45 + 0.2 + 0.8 + 0.9 + 0.8 + 0.6 + 0.9 + 0.6 = 5.45
    \]
    
    \item 
    \[
    E(2X + 3Y - 1) = 2E(X) + 3E(Y) - 1 = 2(1.95) + 3(2.8) - 1 = 3.9 + 8.4 - 1 = 11.3
    \]
    
    \item 
    \[
    E(X^2) = \sum_{x} x^2 \cdot p_X(x) = 1^2(0.3) + 2^2(0.45) + 3^2(0.25) = 0.3 + 1.8 + 2.25 = 4.35
    \]
\end{enumerate}
\end{solution}
\end{eg}

\begin{exercise}[DIY: Expectation of Product for Independent RVs]\index{expectation!product|see{expectation}}
Prove that if $X$ and $Y$ are independent discrete random variables, then:
\[
E(XY) = E(X)E(Y)
\]

\textbf{Hint:} Use the fact that $p_{X,Y}(x, y) = p_X(x) \cdot p_Y(y)$ when $X$ and $Y$ are independent.
\end{exercise}

\begin{proof}
\[
E(XY) = \sum_{x} \sum_{y} xy \cdot p_{X,Y}(x, y) = \sum_{x} \sum_{y} xy \cdot p_X(x) \cdot p_Y(y)
\]
\[
= \sum_{x} x \cdot p_X(x) \sum_{y} y \cdot p_Y(y) = E(X) \cdot E(Y)
\]
\end{proof}

\begin{exercise}[DIY: Law of Total Expectation]\index{law of total expectation}
Prove the law of total expectation for discrete random variables:
\[
E(Y) = \sum_{x} E(Y \mid X = x) \cdot p_X(x)
\]
where $E(Y \mid X = x) = \sum_{y} y \cdot p_{Y|X}(y \mid x)$ is the conditional expectation.
\end{exercise}

\begin{proof}
\[
\sum_{x} E(Y \mid X = x) \cdot p_X(x) = \sum_{x} \left[ \sum_{y} y \cdot p_{Y|X}(y \mid x) \right] \cdot p_X(x)
\]
\[
= \sum_{x} \sum_{y} y \cdot \frac{p_{X,Y}(x, y)}{p_X(x)} \cdot p_X(x) = \sum_{x} \sum_{y} y \cdot p_{X,Y}(x, y) = E(Y)
\]
\end{proof}

\begin{exercise}[DIY: Expectation of Function of Independent RVs]
Prove that if $X$ and $Y$ are independent discrete random variables and $g: \R \to \R$ and $h: \R \to \R$ are functions, then:
\[
E[g(X)h(Y)] = E[g(X)] \cdot E[h(Y)]
\]

\textbf{Hint:} This generalizes the result $E(XY) = E(X)E(Y)$ for independent random variables.
\end{exercise}

\begin{proof}
\[
E[g(X)h(Y)] = \sum_{x} \sum_{y} g(x)h(y) \cdot p_{X,Y}(x, y) = \sum_{x} \sum_{y} g(x)h(y) \cdot p_X(x) \cdot p_Y(y)
\]
\[
= \sum_{x} g(x) \cdot p_X(x) \sum_{y} h(y) \cdot p_Y(y) = E[g(X)] \cdot E[h(Y)]
\]
\end{proof}

\newpage

\section{Exam-Style Practice Problems}
\index{practice problems}

\begin{exercise}
Let $X$ and $Y$ be discrete random variables with joint PMF:

\begin{center}
\renewcommand{\arraystretch}{1.5}
\begin{tabular}{c|ccc}
$p_{X,Y}(x,y)$ & $X=1$ & $X=2$ & $X=3$ \\ \hline
$Y=2$ & 0.1 & 0.15 & 0.05 \\
$Y=3$ & 0.2 & 0.15 & 0.1 \\
$Y=4$ & 0.1 & 0.05 & 0.1 \\
\end{tabular}
\end{center}

\begin{enumerate}[label=(\alph*)]
    \item Find the marginal PMFs of $X$ and $Y$.
    \item Find $P(X = 2 \mid Y = 3)$.
    \item Are $X$ and $Y$ independent? Justify your answer.
    \item Find $E(X)$, $E(Y)$, and $E(XY)$.
    \item \textbf{DIY:} Verify that $E(XY) = E(X)E(Y)$ using the definition of expectation. Does this contradict independence? Explain.
\end{enumerate}
\end{exercise}

\begin{proof}
\begin{enumerate}[label=(\alph*)]
    \item 
    \[
    p_X(1) = 0.1 + 0.2 + 0.1 = 0.4, \quad p_X(2) = 0.15 + 0.15 + 0.05 = 0.35, \quad p_X(3) = 0.05 + 0.1 + 0.1 = 0.25
    \]
    \[
    p_Y(2) = 0.1 + 0.15 + 0.05 = 0.3, \quad p_Y(3) = 0.2 + 0.15 + 0.1 = 0.45, \quad p_Y(4) = 0.1 + 0.05 + 0.1 = 0.25
    \]
    
    \item 
    \[
    P(X = 2 \mid Y = 3) = \frac{0.15}{0.45} = \frac{1}{3}
    \]
    
    \item Check: $p_{X,Y}(1, 2) = 0.1$ and $p_X(1) \cdot p_Y(2) = 0.4 \cdot 0.3 = 0.12$. Since $0.1 \neq 0.12$, $X$ and $Y$ are \textbf{not independent}.
    
    \item 
    \[
    E(X) = 1(0.4) + 2(0.35) + 3(0.25) = 0.4 + 0.7 + 0.75 = 1.85
    \]
    \[
    E(Y) = 2(0.3) + 3(0.45) + 4(0.25) = 0.6 + 1.35 + 1.0 = 2.95
    \]
    \[
    E(XY) = 1 \cdot 2(0.1) + 1 \cdot 3(0.2) + 1 \cdot 4(0.1) + 2 \cdot 2(0.15) + 2 \cdot 3(0.15)
    \]
    \[
    + 2 \cdot 4(0.05) + 3 \cdot 2(0.05) + 3 \cdot 3(0.1) + 3 \cdot 4(0.1)
    \]
    \[
    = 0.2 + 0.6 + 0.4 + 0.6 + 0.9 + 0.4 + 0.3 + 0.9 + 1.2 = 5.5
    \]
    
    \item \textbf{DIY Solution:} We already calculated $E(X) = 1.85$, $E(Y) = 2.95$, and $E(XY) = 5.5$.
    
    Now $E(X)E(Y) = 1.85 \cdot 2.95 = 5.4575 \neq 5.5 = E(XY)$.
    
    This does not contradict independence. In fact, since $E(XY) \neq E(X)E(Y)$, this confirms that $X$ and $Y$ are \textbf{not independent}. If they were independent, we would have $E(XY) = E(X)E(Y)$.
\end{enumerate}
\end{proof}

\begin{exercise}
Let $X$ and $Y$ be independent discrete random variables with:
\[
p_X(1) = 0.3, \quad p_X(2) = 0.5, \quad p_X(3) = 0.2
\]
\[
p_Y(2) = 0.4, \quad p_Y(3) = 0.3, \quad p_Y(4) = 0.2, \quad p_Y(5) = 0.1
\]

\begin{enumerate}[label=(\alph*)]
    \item Find the joint PMF $p_{X,Y}(x, y)$.
    \item Find $E(X)$, $E(Y)$, and $E(XY)$.
    \item Find $E(2X - 3Y + 1)$.
\end{enumerate}
\end{exercise}

\begin{proof}
\begin{enumerate}[label=(\alph*)]
    \item Since $X$ and $Y$ are independent, $p_{X,Y}(x, y) = p_X(x) \cdot p_Y(y)$:
    
    \begin{center}
    \renewcommand{\arraystretch}{1.5}
    \begin{tabular}{c|ccc}
    $p_{X,Y}(x,y)$ & $X=1$ & $X=2$ & $X=3$ \\ \hline
    $Y=2$ & 0.12 & 0.20 & 0.08 \\
    $Y=3$ & 0.09 & 0.15 & 0.06 \\
    $Y=4$ & 0.06 & 0.10 & 0.04 \\
    $Y=5$ & 0.03 & 0.05 & 0.02 \\
    \end{tabular}
    \end{center}
    
    \item 
    \[
    E(X) = 1(0.3) + 2(0.5) + 3(0.2) = 0.3 + 1.0 + 0.6 = 1.9
    \]
    \[
    E(Y) = 2(0.4) + 3(0.3) + 4(0.2) + 5(0.1) = 0.8 + 0.9 + 0.8 + 0.5 = 3.0
    \]
    
    Since $X$ and $Y$ are independent, $E(XY) = E(X)E(Y) = 1.9 \cdot 3.0 = 5.7$.
    
    \item 
    \[
    E(2X - 3Y + 1) = 2E(X) - 3E(Y) + 1 = 2(1.9) - 3(3.0) + 1 = 3.8 - 9.0 + 1 = -4.2
    \]
\end{enumerate}
\end{proof}

\begin{exercise}
Let $X$ and $Y$ have joint PMF:

\begin{center}
\renewcommand{\arraystretch}{1.5}
\begin{tabular}{c|cc}
$p_{X,Y}(x,y)$ & $X=1$ & $X=2$ \\ \hline
$Y=2$ & $a$ & $b$ \\
$Y=3$ & $c$ & $d$ \\
\end{tabular}
\end{center}

where $a + b + c + d = 1$ and all probabilities are non-negative.

\begin{enumerate}[label=(\alph*)]
    \item Express the marginal PMFs in terms of $a, b, c, d$.
    \item Find conditions on $a, b, c, d$ such that $X$ and $Y$ are independent.
    \item If $E(X) = 1.6$ and $E(Y) = 2.4$, find $a, b, c, d$ assuming independence.
    \item \textbf{DIY:} Show that if $X$ and $Y$ are independent, then the conditions in part (b) reduce to a single constraint. What is this constraint?
\end{enumerate}
\end{exercise}

\begin{proof}
\begin{enumerate}[label=(\alph*)]
    \item 
    \[
    p_X(1) = a + c, \quad p_X(2) = b + d
    \]
    \[
    p_Y(2) = a + b, \quad p_Y(3) = c + d
    \]
    
    \item For independence, we need $p_{X,Y}(x, y) = p_X(x) \cdot p_Y(y)$ for all $x, y$:
    \[
    a = (a+c)(a+b), \quad b = (b+d)(a+b), \quad c = (a+c)(c+d), \quad d = (b+d)(c+d)
    \]
    
    \item If independent and $E(X) = 1.6$, $E(Y) = 2.4$:
    \[
    1 \cdot (a+c) + 2 \cdot (b+d) = 1.6 \implies a + c + 2b + 2d = 1.6
    \]
    \[
    2 \cdot (a+b) + 3 \cdot (c+d) = 2.4 \implies 2a + 2b + 3c + 3d = 2.4
    \]
    
    Also, $a + b + c + d = 1$.
    
    Using independence: $a = (a+c)(a+b)$, $b = (b+d)(a+b)$, $c = (a+c)(c+d)$, $d = (b+d)(c+d)$.
    
    Let $p = a+c$ and $q = a+b$. Then $p_X(1) = p$, $p_X(2) = 1-p$, $p_Y(2) = q$, $p_Y(3) = 1-q$.
    
    From $E(X) = 1.6$: $p + 2(1-p) = 1.6 \implies 2 - p = 1.6 \implies p = 0.4$.
    
    From $E(Y) = 2.4$: $2q + 3(1-q) = 2.4 \implies 3 - q = 2.4 \implies q = 0.6$.
    
    Therefore:
    \[
    a = p \cdot q = 0.4 \cdot 0.6 = 0.24
    \]
    \[
    b = (1-p) \cdot q = 0.6 \cdot 0.6 = 0.36
    \]
    \[
    c = p \cdot (1-q) = 0.4 \cdot 0.4 = 0.16
    \]
    \[
    d = (1-p) \cdot (1-q) = 0.6 \cdot 0.4 = 0.24
    \]
    
    Verification: $0.24 + 0.36 + 0.16 + 0.24 = 1.00$
    
    \item \textbf{DIY Solution:} If $X$ and $Y$ are independent, then $p_{X,Y}(x, y) = p_X(x) \cdot p_Y(y)$ for all $x, y$.
    
    Let $p = p_X(1) = a + c$ and $q = p_Y(2) = a + b$. Then:
    \[
    a = p \cdot q, \quad b = (1-p) \cdot q, \quad c = p \cdot (1-q), \quad d = (1-p) \cdot (1-q)
    \]
    
    The single constraint is that $a + b + c + d = 1$, which is automatically satisfied:
    \[
    pq + (1-p)q + p(1-q) + (1-p)(1-q) = q + (1-q) = 1
    \]
    
    So independence reduces the four parameters $(a, b, c, d)$ to just two parameters $(p, q)$.
\end{enumerate}
\end{proof}

\begin{exercise}
Let $X$ and $Y$ be random variables with:
\[
E(X) = 2, \quad E(Y) = 3
\]

Find:
\begin{enumerate}[label=(\alph*)]
    \item $E(3X + 2Y - 5)$
    \item $E(X^2)$ if $\text{Var}(X) = 4$
    \item $E(Y^2)$ if $\text{Var}(Y) = 9$
\end{enumerate}
\end{exercise}

\begin{proof}
\begin{enumerate}[label=(\alph*)]
    \item 
    \[
    E(3X + 2Y - 5) = 3E(X) + 2E(Y) - 5 = 3(2) + 2(3) - 5 = 6 + 6 - 5 = 7
    \]
    
    \item 
    \[
    E(X^2) = \text{Var}(X) + [E(X)]^2 = 4 + 2^2 = 8
    \]
    
    \item 
    \[
    E(Y^2) = \text{Var}(Y) + [E(Y)]^2 = 9 + 3^2 = 18
    \]
\end{enumerate}
\end{proof}

\begin{exercise}
Let $X$ and $Y$ have joint PMF:

\begin{center}
\renewcommand{\arraystretch}{1.5}
\begin{tabular}{c|ccc}
$p_{X,Y}(x,y)$ & $X=1$ & $X=2$ & $X=3$ \\ \hline
$Y=2$ & 0.06 & 0.12 & 0.12 \\
$Y=3$ & 0.08 & 0.16 & 0.16 \\
$Y=4$ & 0.06 & 0.12 & 0.12 \\
\end{tabular}
\end{center}

\begin{enumerate}[label=(\alph*)]
    \item Find the marginal PMFs.
    \item Show that $X$ and $Y$ are independent.
    \item Find $E(XY)$ and verify that $E(XY) = E(X)E(Y)$.
\end{enumerate}
\end{exercise}

\begin{proof}
\begin{enumerate}[label=(\alph*)]
    \item 
    \[
    p_X(1) = 0.20, \quad p_X(2) = 0.40, \quad p_X(3) = 0.40
    \]
    \[
    p_Y(2) = 0.30, \quad p_Y(3) = 0.40, \quad p_Y(4) = 0.30
    \]
    
    \item Check: $p_{X,Y}(1, 2) = 0.06 = p_X(1) \cdot p_Y(2) = 0.20 \cdot 0.30 = 0.06$
    
    $p_{X,Y}(2, 3) = 0.16 = p_X(2) \cdot p_Y(3) = 0.40 \cdot 0.40 = 0.16$
    
    $p_{X,Y}(3, 4) = 0.12 = p_X(3) \cdot p_Y(4) = 0.40 \cdot 0.30 = 0.12$
    
    All checks pass, so $X$ and $Y$ are \textbf{independent}.
    
    \item 
    \[
    E(X) = 1(0.20) + 2(0.40) + 3(0.40) = 0.20 + 0.80 + 1.20 = 2.20
    \]
    \[
    E(Y) = 2(0.30) + 3(0.40) + 4(0.30) = 0.60 + 1.20 + 1.20 = 3.00
    \]
    \[
    E(XY) = \sum_{x} \sum_{y} xy \cdot p_{X,Y}(x, y) = 1 \cdot 2(0.06) + 1 \cdot 3(0.08) + 1 \cdot 4(0.06)
    \]
    \[
    + 2 \cdot 2(0.12) + 2 \cdot 3(0.16) + 2 \cdot 4(0.12) + 3 \cdot 2(0.12) + 3 \cdot 3(0.16) + 3 \cdot 4(0.12) = 6.60
    \]
    
    Verification: $E(X)E(Y) = 2.20 \cdot 3.00 = 6.60$
\end{enumerate}
\end{proof}

\newpage

\section{Summary and Key Formulas}
\index{summary}
\index{formulas}

\subsection{Joint Distribution Formulas}

\begin{center}
\renewcommand{\arraystretch}{1.8}
\begin{tabular}{|l|c|}
\hline
\textbf{Concept} & \textbf{Formula} \\ \hline
\hline
Joint PMF & $p_{X,Y}(x,y) = P(X=x, Y=y)$ \\ \hline
Marginal PMF of $X$ & $p_X(x) = \sum_{y} p_{X,Y}(x,y)$ \\ \hline
Marginal PMF of $Y$ & $p_Y(y) = \sum_{x} p_{X,Y}(x,y)$ \\ \hline
Conditional PMF & $p_{Y|X}(y|x) = \frac{p_{X,Y}(x,y)}{p_X(x)}$ \\ \hline
Independence & $p_{X,Y}(x,y) = p_X(x) \cdot p_Y(y)$ for all $x,y$ \\ \hline
$E[g(X,Y)]$ & $\sum_{x} \sum_{y} g(x,y) \cdot p_{X,Y}(x,y)$ \\ \hline
$E(aX + bY + c)$ & $aE(X) + bE(Y) + c$ (always) \\ \hline
$E(XY)$ & $\sum_{x} \sum_{y} xy \cdot p_{X,Y}(x,y)$ \\ \hline
\end{tabular}
\end{center}

\subsection{Key Properties}

\begin{enumerate}
    \item \textbf{Linearity of Expectation:} $E(aX + bY + c) = aE(X) + bE(Y) + c$ holds \textbf{always}, even if $X$ and $Y$ are dependent.
    
    \item \textbf{Independence:} If $X$ and $Y$ are independent:
    \begin{itemize}
        \item $E(XY) = E(X)E(Y)$
    \end{itemize}
\end{enumerate}

\subsection{Problem-Solving Strategy}

\begin{enumerate}
    \item \textbf{Always start by finding marginals} - they're needed for everything else.
    
    \item \textbf{For independence checks:} Verify $p_{X,Y}(x,y) = p_X(x) \cdot p_Y(y)$ for all $(x,y)$.
    
    \item \textbf{For expectations:} Use linearity whenever possible - it's the easiest approach.
    
    \item \textbf{For conditional probabilities:} Use $P(A \mid B) = \frac{P(A \cap B)}{P(B)}$.
\end{enumerate}

\newpage

\section*{Key Terms and Concepts}
\addcontentsline{toc}{section}{Key Terms and Concepts}

\begin{description}[leftmargin=2cm, style=nextline]
    \item[Joint PMF] \index{joint PMF} The probability mass function $p_{X,Y}(x,y) = P(X=x, Y=y)$ for discrete random variables $X$ and $Y$. It satisfies: (1) $p_{X,Y}(x,y) \geq 0$ for all $x,y$, and (2) $\sum_x \sum_y p_{X,Y}(x,y) = 1$.
    
    \item[Marginal PMF] \index{marginal PMF} The probability mass function of a single random variable obtained by summing the joint PMF over all values of the other variable: $p_X(x) = \sum_y p_{X,Y}(x,y)$ and $p_Y(y) = \sum_x p_{X,Y}(x,y)$.
    
    \item[Conditional PMF] \index{conditional PMF} The probability mass function of one random variable given the value of another: $p_{Y|X}(y|x) = \frac{p_{X,Y}(x,y)}{p_X(x)}$ (provided $p_X(x) > 0$). Conditional PMFs are valid PMFs.
    
    \item[Independence] \index{independence} Two discrete random variables $X$ and $Y$ are independent if $p_{X,Y}(x,y) = p_X(x) \cdot p_Y(y)$ for all $x,y$. Equivalently, $p_{Y|X}(y|x) = p_Y(y)$ for all $x,y$ with $p_X(x) > 0$.
    
    \item[Linearity of Expectation] \index{linearity of expectation} For any constants $a, b, c$: $E(aX + bY + c) = aE(X) + bE(Y) + c$. This holds \textbf{always}, regardless of whether $X$ and $Y$ are independent.
    
    \item[Expectation of Product] \index{expectation!product} For independent random variables: $E(XY) = E(X)E(Y)$. More generally, $E[g(X)h(Y)] = E[g(X)]E[h(Y)]$ for independent $X$ and $Y$. Note: $E(XY) = E(X)E(Y)$ is necessary but not sufficient for independence.
    
    \item[Law of Total Expectation] \index{law of total expectation} $E(Y) = \sum_x E(Y|X=x) \cdot p_X(x)$, where $E(Y|X=x) = \sum_y y \cdot p_{Y|X}(y|x)$ is the conditional expectation.
    
    \item[Problem-Solving Strategy] \index{problem solving} Always start by finding marginals. For independence checks, verify $p_{X,Y}(x,y) = p_X(x) \cdot p_Y(y)$ for all $(x,y)$. For expectations, use linearity whenever possible.
\end{description}

